\section{Evaluasi}
\label{sec:evaluasi-bab2}

\subsection{Soal Pilihan Ganda}

\begin{enumerate}
    \item Prinsip yang menyatakan bahwa instruksi dan data disimpan
          pada memori yang sama disebut \dots
          \begin{enumerate}
              \item \eng{pipelining}
              \item \eng{stored-program}
              \item \eng{memory-mapped I/O}
              \item \eng{little-endian}
          \end{enumerate}

    \item Urutan tahapan siklus instruksi yang benar adalah \dots
          \begin{enumerate}
              \item \eng{decode} $\rightarrow$ \eng{fetch}
                    $\rightarrow$ \eng{execute}
              \item \eng{execute} $\rightarrow$ \eng{fetch}
                    $\rightarrow$ \eng{decode}
              \item \eng{fetch} $\rightarrow$ \eng{decode}
                    $\rightarrow$ \eng{execute}
              \item \eng{fetch} $\rightarrow$ \eng{execute}
                    $\rightarrow$ \eng{decode}
          \end{enumerate}

    \item Nilai \dec{-10} dalam 8-bit \eng{two's complement} adalah
          \dots
          \begin{enumerate}
              \item \hex{0A} \item \hex{0F}
              \item \hex{F5} \item \hex{F6}
          \end{enumerate}

    \item Hasil \asm{ADD AL, 01h} dengan \reg{AL} = \hex{7F} adalah
          \dots
          \begin{enumerate}
              \item \reg{AL} = \hex{80}, \flag{OF} = 0,
                    \flag{CF} = 0
              \item \reg{AL} = \hex{80}, \flag{OF} = 1,
                    \flag{CF} = 0
              \item \reg{AL} = \hex{80}, \flag{OF} = 1,
                    \flag{CF} = 1
              \item \reg{AL} = \hex{00}, \flag{OF} = 0,
                    \flag{CF} = 1
          \end{enumerate}

    \item Sandian ASCII untuk karakter \code{'A'} adalah \dots
          \begin{enumerate}
              \item \hex{0A} \item \hex{41}
              \item \hex{61} \item \hex{45}
          \end{enumerate}
\end{enumerate}

\subsection{Soal Uraian}

\begin{enumerate}
    \item Jelaskan model mesin penyimpan program beserta keempat
          komponennya. Kemudian jelaskan satu konsekuensi penting
          dari prinsip \eng{stored-program} bagi cara program dimuat
          dan dijalankan.

    \item Jelaskan perbedaan antara \eng{carry} dan \eng{overflow}.
          Sertakan satu contoh operasi yang menghasilkan \flag{CF} = 1
          tetapi \flag{OF} = 0, dan satu contoh yang menghasilkan
          \flag{CF} = 0 tetapi \flag{OF} = 1.

    \item Sebuah variabel berisi \hex{32}. Jelaskan dua tafsiran yang
          mungkin atas isi variabel tersebut, dan sebutkan konteks
          yang menentukan tafsiran mana yang benar.

    \item Berdasarkan Gambar~\ref{fig:peta-konsep}, susun rencana
          belajar Anda untuk enam minggu pertama. Jelaskan mengapa
          urutan bab tidak dapat ditukar sembarangan.
\end{enumerate}

\subsection{Kisi-kisi Penilaian}

\begin{table}[H]
    \centering
    \caption{Kisi-kisi evaluasi Bab~\ref{chap:landasan-teori}.}
    \label{tab:kisi-bab2}
    \small
    \begin{tabular}{@{}p{5.4cm}p{5.2cm}c@{}}
        \toprule
        \textbf{Aspek yang diukur} & \textbf{Bentuk soal} &
        \textbf{Bobot} \\
        \midrule
        Konsep model mesin dan siklus instruksi & Pilihan ganda 1, 2;
        Uraian 1 & 25\% \\
        \addlinespace
        Representasi bilangan bertanda & Pilihan ganda 3; Uraian 2 &
        25\% \\
        \addlinespace
        Pembedaan \eng{carry} dan \eng{overflow} & Pilihan ganda 4;
        Uraian 2 & 20\% \\
        \addlinespace
        Penyandian data dan tafsirannya & Pilihan ganda 5; Uraian 3 &
        15\% \\
        \addlinespace
        Kemampuan membaca peta konsep buku & Uraian 4 & 15\% \\
        \bottomrule
    \end{tabular}
\end{table}

\begin{checklistbox}
\textbf{Checklist Pencapaian Fondasi Bab~\ref{chap:landasan-teori}}

Mahasiswa dinyatakan siap memasuki bab materi apabila:
\begin{itemize}
    \item[\checked] dapat menjelaskan model von Neumann dan siklus
          instruksi;
    \item[\checked] dapat mengonversi bilangan biner, heksadesimal,
          dan desimal dengan benar tanpa alat bantu;
    \item[\checked] dapat menyatakan bilangan negatif dalam
          \eng{two's complement} dan memverifikasi hasilnya;
    \item[\checked] dapat membedakan \eng{carry} dan \eng{overflow}
          beserta flag yang menandainya;
    \item[\checked] dapat menjelaskan perbedaan nilai biner dan
          sandian karakter;
    \item[\checked] dapat menjelaskan urutan ketergantungan antarbab
          melalui peta konsep.
\end{itemize}

\vspace{0.4em}
Seluruh butir di atas bersifat prasyarat untuk
Bab~\ref{chap:pengantar-assembly}.
\end{checklistbox}
